\chapter*{INTRODUCTION GÉNÉRALE}
\addcontentsline{toc}{chapter}{INTRODUCTION GÉNÉRALE}

% L'introduction est volontairement non numérotée : aucune section 0.x ne doit apparaître.
\setcounter{secnumdepth}{-1}

\section{Contexte général}

Les journaux produits par les systèmes informatiques constituent une mémoire opérationnelle de l'infrastructure. Ils enregistrent les connexions, les erreurs, les changements de configuration, les accès aux services, les requêtes applicatives et les alertes de sécurité. Dans un environnement de cybersécurité, ces traces servent à comprendre ce qui s'est produit, à repérer des comportements inhabituels et à documenter les décisions prises pendant l'analyse. Le guide NIST consacré à la gestion des journaux rappelle d'ailleurs que leur valeur dépend autant de la collecte et de la conservation que de l'analyse proprement dite \cite{nist80092}.

La difficulté ne vient pas seulement du volume. Les sources sont également hétérogènes par leur format, leur granularité et leur signification. Un événement Windows n'a pas la même structure qu'une ligne syslog Linux. Une requête Apache décrit une interaction applicative, tandis qu'un export Wazuh contient déjà des règles, des niveaux de sévérité et des identifiants d'agent. Les jeux de données réseau ajoutent une autre perspective : ils représentent des flux par des variables numériques et par des étiquettes de trafic. Cette diversité est utile pour la supervision, mais elle empêche une analyse uniforme sans étape préalable de structuration.

Dans une petite infrastructure ou dans un contexte universitaire, l'analyste doit souvent passer d'un outil à l'autre, rechercher les champs pertinents et comparer des alertes qui ne parlent pas le même langage. L'analyse manuelle devient alors lente et difficile à reproduire. Elle dépend de la connaissance de la source, de la qualité des filtres et de la capacité à conserver le contexte d'un événement. L'automatisation peut réduire cette charge, à condition de ne pas transformer les journaux en simples scores impossibles à interpréter.

Les travaux sur la détection d'intrusion ont montré l'intérêt des traces d'audit pour caractériser les déviations d'un comportement attendu \cite{denning1987intrusion,axelsson2000intrusion}. La détection d'anomalies et l'apprentissage automatique ont ensuite élargi les possibilités de traitement, notamment lorsque les règles connues ne suffisent plus \cite{chandola2009anomaly,buczak2016survey}. Ces méthodes restent toutefois sensibles aux labels, au déséquilibre des classes, à la dérive des distributions et au protocole de séparation des données \cite{khraisat2019survey,sommer2010outside}. Une chaîne d'analyse crédible doit donc relier le modèle à ses données, à son protocole et à la décision humaine qui suit le signal.

\section{Constat scientifique et technique}

Trois constats structurent le problème étudié. Premièrement, les données disponibles ne sont pas homogènes. Certaines sources sont des journaux bruts, d'autres des exports enrichis et d'autres encore des tableaux déjà étiquetés. Un même modèle ne peut pas être appliqué sans précaution à ces familles différentes. Deuxièmement, la performance d'un détecteur ne suffit pas à établir son utilité. Un score élevé sur une partition aléatoire peut masquer une dépendance forte au fichier ou au scénario d'origine. Troisièmement, l'exécution elle-même doit être observable : une tâche perdue, un résultat non acquitté ou une décision non tracée compromet la reproductibilité de l'analyse.

Ces constats déplacent la question au-delà du choix d'un algorithme. Il faut construire une chaîne capable de collecter les sources, d'extraire des champs communs, de conserver les particularités, de choisir un traitement compatible avec chaque famille et de présenter le résultat dans un contexte compréhensible. Le système doit également indiquer ce qu'il ne sait pas démontrer. Une anomalie produite sans label fiable reste un candidat ; une campagne distribuée réalisée dans VirtualBox reste une preuve de laboratoire ; une mise à jour de modèle exécutée sur un cas contrôlé ne constitue pas encore un apprentissage continu en production.

Le prototype Ariel Logminer est conçu dans cet espace d'ingénierie. Il associe des collecteurs, des parseurs, une normalisation, un routeur par famille, des modèles supervisés et non supervisés, une corrélation, une mémoire d'exécution, un audit et un dashboard. L'ambition est d'intégrer ces fonctions dans une architecture locale mesurable, et non de proposer un nouvel algorithme de détection ou un SOC industriel complet.

\section{Problématique}

Les journaux provenant de plusieurs systèmes et sources sont hétérogènes, volumineux et difficiles à traiter de manière uniforme. Les approches existantes couvrent chacune une partie du problème : un collecteur facilite l'acquisition, un parseur reconnaît un format, un modèle calcule un score et un outil de visualisation affiche des événements. La difficulté apparaît lorsqu'il faut relier ces fonctions sans perdre l'origine de la donnée ni le contexte de la décision.

La problématique de ce mémoire porte donc sur l'intégration cohérente de la collecte, du parsing, de la normalisation, du routage, de la détection, de la corrélation, de la mémoire, de l'audit et de la visualisation. Cette intégration doit rester suffisamment modulaire pour accueillir plusieurs familles de journaux et suffisamment explicable pour qu'un analyste puisse remonter du signal vers la source et le protocole qui l'ont produit.

\begin{quote}
Comment concevoir un prototype autonome, distribué au sens logiciel, modulaire et exploitable par un analyste, capable de détecter et de prioriser des anomalies candidates dans des journaux systèmes et réseaux hétérogènes par traitements périodiques, par micro-lots et par agents multi-tâches ?
\end{quote}

Cette formulation contient plusieurs qualifications importantes. « Autonome » désigne ici une orchestration logicielle fondée sur des capacités, des règles et des états explicites. « Distribué » renvoie à la séparation en processus et à la validation inter-machines réalisée en laboratoire. Enfin, « anomalie candidate » ne signifie pas attaque confirmée : la décision de sécurité reste soumise au contexte, aux preuves disponibles et à l'analyste.

\section{Question centrale et questions spécifiques de recherche}

La question centrale est étudiée au moyen de six questions opérationnelles. Elles couvrent le schéma de données, l'organisation logicielle, la sélection des modèles, l'interface, la mémoire et l'évaluation.

\begin{enumerate}
  \item \textbf{QR1.} Comment construire un schéma de normalisation rapprochant journaux système, réseau et SIEM sans perdre l'information nécessaire ?
  \item \textbf{QR2.} Comment organiser des agents spécialisés dans une architecture modulaire, auditable et extensible ?
  \item \textbf{QR3.} Comment sélectionner automatiquement un modèle adapté à la famille traitée ?
  \item \textbf{QR4.} Comment présenter les anomalies dans un tableau de bord permettant compréhension, filtrage et validation ?
  \item \textbf{QR5.} Comment exploiter une mémoire persistante pour améliorer la priorisation et organiser la mise à jour contrôlée ?
  \item \textbf{QR6.} Comment évaluer le prototype par métriques, captures, latence et ressources ?
\end{enumerate}

Ces questions ne sont pas traitées comme six performances indépendantes. Elles forment une chaîne : la normalisation conditionne le routage, le routage dépend de l'organisation des agents, les sorties des modèles doivent être présentées et auditées, et l'ensemble doit être évalué avec des protocoles qui rendent les résultats interprétables.

\section{Hypothèses de recherche}

Les hypothèses suivantes sont formulées comme des propositions à examiner expérimentalement. Leur statut final est discuté plus loin dans le mémoire ; aucune n'est déclarée validée dans cette introduction.

\begin{enumerate}
  \item \textbf{H1.} Des journaux hétérogènes peuvent être rapprochés par un noyau commun sans supprimer les informations originales utiles à l'audit.
  \item \textbf{H2.} Des agents spécialisés peuvent améliorer la modularité, la traçabilité et la robustesse par rapport à une chaîne monolithique.
  \item \textbf{H3.} Le routage par famille peut limiter les incompatibilités entre données, caractéristiques et modèles.
  \item \textbf{H4.} Des modèles légers conservent un intérêt dans une chaîne locale reproductible avec métriques, corrélation et validation humaine.
  \item \textbf{H5.} Une mémoire persistante peut soutenir l'adaptation progressive, la priorisation et la mise à jour contrôlée.
\end{enumerate}

La formulation de ces hypothèses impose une lecture prudente des résultats. Une architecture peut être modulaire sans offrir un gain de débit. Un routeur peut fonctionner correctement sans produire un gain prédictif systématique. Une mémoire peut conserver des décisions sans constituer un mécanisme d'apprentissage continu. Ces distinctions guideront l'analyse expérimentale et empêcheront de confondre fonctionnement, efficacité et généralisation.

\section{Objectif général}

L'objectif général est de concevoir et d'évaluer Ariel Logminer, un prototype local et modulaire capable d'organiser l'analyse d'anomalies candidates dans des journaux systèmes et réseaux hétérogènes. Le prototype doit relier l'acquisition des données à la présentation des résultats en passant par la normalisation, le routage, la détection, la corrélation, la mémoire et l'audit.

\section{Objectifs spécifiques}

Les objectifs spécifiques découlent directement des questions de recherche :

\begin{enumerate}
  \item caractériser les familles de journaux et de flux retenues pour l'évaluation ;
  \item définir un schéma commun qui conserve les champs utiles et l'origine des sources ;
  \item concevoir une architecture d'agents spécialisés avec contrats, capacités, mémoire et mécanismes d'audit ;
  \item implémenter les collecteurs, parseurs, normaliseurs, routeurs, détecteurs, corrélateurs et services de visualisation ;
  \item comparer des modèles légers avec des séparations aléatoires et des holdouts par fichier ou scénario ;
  \item évaluer séparément les traitements HDFS/BGL, la robustesse multiformat, le routage et la corrélation contrôlée ;
  \item mesurer le comportement opérationnel au moyen de campagnes de ressources, de latence, de reprise et de distribution de laboratoire ;
  \item documenter les limites de validité, la traçabilité des données et les conditions de reproduction.
\end{enumerate}

\section{Démarche méthodologique}

La démarche commence par une analyse des formats de journaux et des exigences d'une chaîne d'analyse exploitable. La conception sépare ensuite les responsabilités : les adaptateurs acquièrent les sources, les parseurs extraient les événements, le normaliseur construit un noyau commun et le routeur sélectionne une famille de traitement. Les détecteurs produisent des signaux qui sont ensuite enrichis par la corrélation, l'audit et le dashboard.

Le prototype est évalué sur quatre catégories de données. Les datasets publics servent aux comparaisons supervisées ou aux séquences de logs lorsqu'une référence et des labels sont disponibles. Les journaux collectés localement, tels que les événements Windows, apportent une observation de formats réellement produits sur la machine d'expérimentation. Les exports locaux, notamment Wazuh, sont traités comme des données déjà enrichies par un outil de sécurité. Enfin, des entrées synthétiques contrôlées servent à isoler des mécanismes précis comme la corrélation et la reprise après panne.

Les expériences finales comprennent une comparaison CICIDS2017 entre partition aléatoire et holdouts par scénario, une comparaison de plusieurs modèles sur des graines appariées, une validation HDFS/BGL avec séparation entraînement--validation--test et Drain3 ajusté uniquement sur l'entraînement, ainsi qu'une validation multiformat. Elles comprennent également l'évaluation du routeur réel, la mise à jour contrôlée des modèles, la corrélation synthétique, les statistiques transversales et les mesures opérationnelles. Les campagnes d'agents et de multi-VM sont présentées comme des validations de laboratoire ; elles ne sont pas extrapolées à une infrastructure industrielle.

L'analyse ne repose pas sur une métrique unique. Les résultats supervisés sont décrits par la précision, le rappel, le F1 et, lorsque le protocole le permet, des métriques complémentaires. Les sorties non supervisées restent des anomalies candidates. Les mesures de latence, de ressources, de débit et de reprise documentent le coût opérationnel. Les hashes, manifestes, scripts et tableaux de preuve relient enfin chaque affirmation à un artefact identifiable.

\section{Délimitations du travail}

Le travail porte sur un prototype exécuté dans un environnement local et contrôlé. Les datasets et journaux utilisés couvrent plusieurs familles, mais ils ne représentent pas toutes les infrastructures ni toutes les distributions de trafic. Certaines copies locales de sources publiques sont vérifiées directement, tandis que d'autres restent documentées avec un niveau de traçabilité plus limité ; cette distinction est conservée dans les manifestes et n'est pas masquée par une appellation générique.

La validation distribuée comprend une campagne locale multiprocessus et une validation inter-machines sous VirtualBox avec Debian et Ubuntu. Cette dernière démontre la consommation d'un flux Redis commun et des scénarios de reprise en laboratoire. Elle ne constitue ni une preuve de haute disponibilité industrielle, ni une validation multi-site sécurisée, ni une mesure de scalabilité à grande échelle.

La mémoire persistante conserve des historiques d'exécution, des erreurs, des décisions et des éléments de feedback. La mise à jour contrôlée peut comparer, promouvoir ou rejeter des modèles dans les scénarios prévus. Ces mécanismes ne doivent pas être décrits comme un apprentissage continu autonome modifiant sans supervision les modèles en production. De même, le dashboard fait l'objet d'une validation fonctionnelle ; aucune étude utilisateur ne permet de conclure à une amélioration mesurée de la compréhension ou de la charge cognitive.

\section{Contributions attendues}

La première contribution est architecturale : elle propose une organisation intégrée reliant acquisition, normalisation, routage, détection, corrélation, audit et visualisation. La deuxième est méthodologique : elle associe des protocoles de séparation, des graines, des métriques et des artefacts de preuve afin de distinguer performance exploratoire et comportement sous holdout plus strict.

La troisième contribution est logicielle. Le dépôt rassemble des collecteurs, des adaptateurs, des agents, des handlers, un routeur, des services locaux, une mémoire et un dashboard. Les contrats et les journaux d'audit rendent les échanges inspectables, tandis que Redis Streams permet d'expérimenter une répartition de tâches et une reprise contrôlée.

La quatrième contribution est expérimentale. Elle met en regard plusieurs familles de données, des modèles légers, des séquences HDFS/BGL, des scénarios multiformat et des mesures d'exécution. Les résultats négatifs, les limites de provenance et l'absence de gain systématique du routage sont conservés, car ils précisent le domaine de validité du prototype.

La cinquième contribution est opérationnelle : elle fournit une chaîne qui peut être examinée par un analyste, depuis la source jusqu'à la décision, avec des vues de filtrage, des détails d'incident, une mémoire d'audit et des preuves reproductibles. Cette contribution reste celle d'une intégration locale ; elle ne revendique ni un nouvel algorithme d'apprentissage ni une plateforme SOC industrielle.

\section{Organisation du mémoire}

Après cette introduction générale, le premier chapitre présente un état de l'art critique sur les journaux, le parsing, la détection d'anomalies, les agents et les systèmes distribués. Le deuxième chapitre expose le modèle proposé et l'architecture d'Ariel Logminer : responsabilités des agents, flux de données, normalisation, routage, communication, mémoire, audit et limites de conception.

Le troisième chapitre décrit l'implémentation du prototype, ses modules, ses interfaces, ses parseurs, ses services et son dashboard. Le quatrième chapitre présente la méthodologie expérimentale et les résultats retenus : données, splits, modèles, Drain3, ressources, latence, robustesse, mise à jour contrôlée, routage, corrélation et campagnes d'agents.

Le cinquième chapitre approfondit la validité scientifique et analyse les limites des protocoles. Le sixième propose une discussion générale des résultats, de la portée de l'architecture et des perspectives. Le septième décrit la reproductibilité, le déploiement local, la validation multi-VM de laboratoire, les précautions d'exploitation et la structure des artefacts. La conclusion générale revient sur les contributions, les hypothèses et les prolongements possibles. Les annexes regroupent les tableaux, manifestes, commandes, captures et éléments de preuve nécessaires à la vérification.

\setcounter{secnumdepth}{3}
